\documentclass[11pt]{article}
\usepackage[utf8]{inputenc}
\usepackage{amsmath,amssymb,amsthm}
\usepackage[margin=1in]{geometry}
\usepackage{hyperref}
\usepackage{xcolor}
\usepackage{enumitem}
\usepackage{newunicodechar}
\newunicodechar{ℝ}{\ensuremath{\mathbb{R}}}
\newunicodechar{ℕ}{\ensuremath{\mathbb{N}}}

\title{Tide retrospective:\\\emph{Integrability for the generic-$k$ monomial template (A1 follow-up)}}
\author{Daniel Murfet}
\date{18 May 2026}

\begin{document}
\maketitle

\begin{abstract}
A1's retrospective deferred the integrability layer
(\texttt{kth\_integrable\_pow*}) on the grounds that the Gaussian-
domination argument needed case analysis on $k = 1$ vs $k \ge 2$.
This Tide closes the layer with a uniform argument that avoids the
case split: the key bound $x^2 \le 1 + x^{2k}$ for $k \ge 1$ is
immediate (case-split on $|x| \le 1$ vs $|x| \ge 1$, the latter via
\texttt{le\_self\_pow$_0$}). Then refactors the quartic and sextic
integrability lemmas as $k = 2$ / $k = 3$ specialisations. Net
$-17$ lines across the three files; build clean (8298 jobs),
zero \texttt{sorry} / \texttt{axiom} / \texttt{native\_decide}.
\end{abstract}

\section{Setting}

A1 (generic even-monomial 1D template) landed nine theorems
parametrised by $k \ge 1$: the moments, the partition function, and
the expected values. The integrability witnesses\footnote{Named
\texttt{kth\_integrable\_pow} and
\texttt{kth\_integrable\_pow\_pot}.} were deferred to a follow-up
tide. A1's retrospective worried that the seabed's Gaussian-
domination argument relied on $k$-specific square-completion
identities ($k = 2$: $-6/t$; $k = 3$: $-\sqrt{240/t} \cdot 2/3$); a
uniform formula at generic $k$ would require a case-split or a
parametric \texttt{Real.rpow} constant.

In fact a simpler approach exists, recorded here.

\section{The thing formalised}

A private bound:
\begin{itemize}[leftmargin=1.5em]
\item \texttt{sq\_le\_one\_add\_pow\_two\_mul} (private):
  for $k \ge 1$ and $x : \mathbb R$, $x^2 \le 1 + x^{2k}$.
\end{itemize}

Two public theorems:
\begin{itemize}[leftmargin=1.5em]
\item \texttt{kth\_integrable\_pow}: for $k \ge 1$, $n : \mathbb N$,
  $t > 0$, the function $x^n \cdot \exp(-tx^{2k}/(2k)!)$ is Lebesgue
  integrable on $\mathbb R$.
\item \texttt{kth\_integrable\_pow\_pot}: same in
  \texttt{kthPotential} form.
\end{itemize}

Plus four refactors in the quartic and sextic files: the four
\texttt{\_integrable\_pow*} theorems become thin specialisations.

\section{Proof strategy}

The bound $x^2 \le 1 + x^{2k}$ for $k \ge 1$ uses an unconditional
case split on $|x|$:

\begin{itemize}[leftmargin=1.5em]
\item \emph{$|x| \le 1$}: $x^2 \le 1 \le 1 + x^{2k}$ since
  $x^{2k} = (x^2)^k \ge 0$.
\item \emph{$|x| \ge 1$}: $x^{2k} = (x^2)^k \ge x^2$ by
  \texttt{le\_self\_pow$_0$} applied to $x^2 \ge 1$ and $k \ge 1$.
  Then $1 + x^{2k} \ge x^2$.
\end{itemize}

From the bound, $t x^{2k}/(2k)! - (t/(2k)!) x^2 \ge -t/(2k)!$, hence
$\exp(-tx^{2k}/(2k)!) \le \exp(t/(2k)!) \cdot \exp(-(t/(2k)!) x^2)$.
The dominator $|x|^n \cdot \exp(-(t/(2k)!) x^2)$ is Gaussian-times-
poly, integrable by Mathlib's
\texttt{integrable\_rpow\_mul\_exp\_neg\_mul\_sq} with parameter
$b = t/(2k)!$. \texttt{Integrable.mono} closes.

\section{Roadblocks and resolutions}

Three small fixes during the proof:

\begin{enumerate}[leftmargin=1.5em]
\item Missing import: the Mathlib Gaussian-integral file was needed
  for the dominator's integrability.\footnote{Specifically
  \texttt{Mathlib.Analysis.SpecialFunctions.Gaussian.GaussianIntegral}
  for \texttt{integrable\_rpow\_mul\_exp\_neg\_mul\_sq}.}
\item \texttt{le\_self\_pow} renamed to
  \texttt{le\_self\_pow$_0$} in the current Mathlib snapshot.
\item \texttt{nlinarith} couldn't close the bound directly because of
  the implicit identity $t x^{2k} / \mathrm{fac} = (t/\mathrm{fac})
  \cdot x^{2k}$; added an explicit \texttt{field\_simp}-based
  rewrite.
\end{enumerate}

Each was a small local fix; none required deliberation.

\section{What was Mathlib, what was new}

\paragraph{From Mathlib.} The Gaussian-times-poly dominator,
\texttt{le\_self\_pow$_0$} for the $(x^2)^k \ge x^2$ step,
\texttt{Integrable.mono} and \texttt{.const\_mul}, plus the
standard \texttt{Real.exp} and absolute-value bookkeeping.

\paragraph{From the seabed.} A1's \texttt{kthPotential} definition.
The quartic and sextic integrability lemmas use the new generic
theorems as specialisations.

\paragraph{New here.} The private bound
$x^2 \le 1 + x^{2k}$ and the two integrability theorems, plus the
four specialisations in quartic / sextic.

\section{Lessons}

\paragraph{Case-split on $k$ is not always necessary for $k$-parametric
arguments.} A1's retrospective expected the integrability proof to
require a case-split on $k = 1$ vs $k \ge 2$ because the seabed's
quartic and sextic proofs use $k$-specific square-completion
identities. In fact the simpler bound $x^2 \le 1 + x^{2k}$ works
uniformly --- it just needed to be found. The lesson: when a follow-
up tide looks blocked by a case-split, search for a coarser uniform
bound that avoids the case-split entirely.

\paragraph{\texttt{nlinarith} sees Real.rpow-style arithmetic
better than \texttt{Real.rpow}-styled goals.} The first build attempt
had \texttt{nlinarith} fail to close a step that turned on the
identity $t x^{2k}/fac = (t/fac) \cdot x^{2k}$. An explicit
\texttt{field\_simp}-rewrite to put the goal in the latter form made
the rest of \texttt{nlinarith} routine. Worth recording.

\paragraph{Mathlib renaming.} \texttt{le\_self\_pow} is now
\texttt{le\_self\_pow$_0$} in the current Mathlib snapshot. Other
\texttt{$_0$}-suffixed renames exist (\texttt{pow\_le\_pow\_right$_0$},
etc.) for the \texttt{GroupWithZero}-class versions. Use the
\texttt{$_0$} variants when working with \texttt{ℝ}-valued bounds.

\section{Follow-ups}

\begin{itemize}[leftmargin=1.5em]
\item \emph{Optional: tighten the bound.} The bound
  $x^2 \le 1 + x^{2k}$ is coarse; for $k \ge 2$, tighter bounds
  with $k$-dependent constants are available (e.g.\ the
  $((k-1)/k) \cdot ((2k)!/(kt))^{1/(k-1)}$ minimum). For downstream
  consumers that need a tight asymptotic constant, the present bound
  is enough; for purely integrability use, the coarse bound suffices.
\item \emph{Promote to \texttt{lean-common}.} If a third consumer of
  the bound $x^2 \le 1 + x^{2k}$ materialises outside the laplace
  repo, factor it up. Currently one use site.
\item \emph{Multi-index analogues.} 2D / multi-index generalisations
  of the integrability bound would follow the same template; deferred
  until a downstream consumer asks.
\end{itemize}

\section*{Acknowledgements}

No GPT consult was used (A1-integrability is a focused follow-up of
A1, and A1's pre-consult had set the relevant scope decisions). The
uniform Gaussian-domination bound was found by inspection of the
seabed's sextic argument and noting that
$x^2 \le 1 + x^6 \Rightarrow x^2 \le 1 + x^{2k}$ for any $k \ge 1$.

\end{document}
